\chapter{Inflation Markets}\label{ch:m2:inflation-markets}

In the twelve months to June 2022 American consumer prices rose 9.1\%, the
fastest in more than forty years. It should have been the year of the
inflation-protected bond, whose principal grows with the price index. An
investor who bought a ten-year one on the first trading day of 2022 and marked
it on the last saw its principal grow by 7.7\% and its price fall by 21.3\%:
a loss of 15.2\% on the year. Nothing was wrong with the indexation. The bond
protects against inflation relative to its \emph{real} yield, and that yield
went from $-0.97\%$ to $+1.58\%$ as the Federal Reserve raised rates. This
chapter explains what an inflation-linked bond pays, how its yield separates
into a real rate and the market's price of inflation, how inflation is traded
in swaps, and why the calendar of price releases makes an inflation-linked
position's carry, month by month, one of the most predictable numbers in
fixed income.

\section{Index-linked bonds}

\begin{definition}[Index-linked bond, reference index, deflation floor]\label{def:m2:inflation-markets:linker}
An \emph{index-linked bond}\index{index-linked bond} (a \emph{linker}) is a bond
whose principal, and so each coupon, is scaled by the growth of a consumer
price index since the bond's dated date. The scaling uses the
\emph{reference index}\index{reference index}: the price index of a given day
as the bond's rules define it from the published monthly values. The ratio of
the reference index on a day to its value on the dated date is the
\emph{index ratio}. A \emph{deflation floor}\index{deflation floor} is a promise
that the principal repaid at maturity is never less than the original face
value, whatever the index has done.
\end{definition}

Treasury Inflation-Protected Securities, TIPS, were first auctioned by the
United States in January 1997. They are indexed to the consumer price index
for all urban consumers, not seasonally adjusted, published monthly by the
Bureau of Labor Statistics; they pay a fixed coupon every six months on the
adjusted principal, and at maturity the holder receives the adjusted principal
or the original, whichever is larger. The fixed coupon is a \emph{real} rate:
a TIPS with a 0.125\% coupon pays 0.0625\% of its adjusted principal each half
year, so its cash flows keep their purchasing power.

A monthly index is published about two weeks after the month it measures, so no
bond can be indexed to today's prices. The TIPS rule, in the Treasury's
regulations, sets the \omterm{def:m2:inflation-markets:linker}{reference index} for the first day of a month equal to
the index of the third preceding month, and interpolates linearly for the
other days:
\begin{equation}\label{eq:m2:inflation-markets:refcpi}
\mathrm{Ref}(d) \;=\; I_{M-3} + \frac{t-1}{D}\,\bigl(I_{M-2} - I_{M-3}\bigr),
\end{equation}
where $d$ is day $t$ of month $M$, $D$ the number of days in $M$ and $I_m$ the
published index of month $m$; the \omterm{def:m2:inflation-markets:linker}{reference index} and the index ratio are
truncated to six decimals and rounded to five. The regulation's own example:
with January 1996 at 154.40 and February at 154.90, the \omterm{def:m2:inflation-markets:linker}{reference index} on 15
April 1996 is $154.40 + \tfrac{14}{30}\times 0.50 = 154.63333$.

\begin{omfigure}
\centering
\begin{tikzpicture}[font=\small,x=1.55cm]
\draw[thick,-{Stealth}] (0,0) -- (6.6,0);
\foreach \m/\l in {0/May,1/Jun,2/Jul,3/Aug,4/Sep,5/Oct,6/Nov} {\draw (\m,0.1) -- (\m,-0.1) node[below] {1 \l};}
\draw[omDef,very thick] (3,1.2) -- (4,1.75);
\node[omDef,align=center,font=\footnotesize] at (3.5,2.25) {reference index\\in August};
\fill[omDef] (3,1.2) circle (1.8pt) node[left,font=\footnotesize] {May index};
\fill[omDef] (4,1.75) circle (1.8pt) node[right,font=\footnotesize] {June index};
\draw[omThm,thick,-{Stealth}] (2.39,-0.9) -- (2.39,-0.08);
\node[omThm,align=center,font=\footnotesize] at (2.39,-1.35) {13 July: June\\index published};
\draw[omProp,thick,decorate,decoration={brace,amplitude=4pt,mirror}] (2.39,-1.95) -- (3,-1.95);
\node[omProp,font=\footnotesize,align=center] at (2.7,-2.45) {the August accrual is\\known before August starts};
\end{tikzpicture}

\omcaption{The \omterm{def:m2:inflation-markets:lag}{indexation lag} of a TIPS. The \omterm{def:m2:inflation-markets:linker}{reference index} moves during
August from the May index to the June index, both already published, so the
growth of the principal over August is known from 13 July, the day the June
index was released in 2022. Schematic.}\label{fig:m2:inflation-markets:lag}
\end{omfigure}

\begin{example}[An index ratio]\label{ex:m2:inflation-markets:ratio}
Take an illustrative ten-year TIPS dated 15 January 2022, with a
0.125\% coupon. Its \omterm{def:m2:inflation-markets:linker}{reference index} on that day
interpolates between the October 2021 index, 276.589, and November's, 277.948:
$276.589 + \tfrac{14}{31}\times 1.359 = 277.20274$. On 15 August 2022 it
interpolates between May 2022, 292.296, and June, 296.311: 294.10922. The index
ratio is $294.10922/277.20274 = 1.06099$, so USD~100 million of face value is
then USD~106.099 million of principal, and the coupon paid on 15 July, at an
index ratio of 1.04814, was USD~65\,509.
\end{example}

\begin{remark}[When the index is not published]\label{rem:m2:inflation-markets:contingency}
A bond indexed to a statistic depends on the statistic being produced. The
Treasury's rule provides for its absence: if a month's index is not reported by
the last day of the following month, the Treasury announces a substitute, the
last reported index grown for one month at its last twelve-month rate,
$I_M = I_{M-1}\,(I_{M-1}/I_{M-13})^{1/12}$, and keeps it even if the true
figure appears later. This happened for the first time in 2025: the Bureau of
Labor Statistics collected no prices from 1 October to 12 November during the
lapse in federal appropriations, and no October index was published. On 26
November the Treasury announced 325.604 for October, which is
$324.800 \times (324.800/315.301)^{1/12}$ from the September indices of 2025
and 2024. Every TIPS cash flow that depends on October 2025 uses that number.
\end{remark}

The other large markets use the same design with different indices. The
United Kingdom's index-linked gilts are indexed to the Retail Prices Index;
those first issued before 2005 use an eight-month lag, those issued since a
three-month lag with an interpolated \omterm{def:m2:inflation-markets:linker}{reference index}, and a gilt repays less
than its face value if the index has fallen: there is no \omterm{def:m2:inflation-markets:linker}{deflation floor}.
France issues two families, one indexed to French consumer prices excluding
tobacco and one, since 2001, to the euro-area harmonised index excluding
tobacco, with the same interpolation between the indices of three and two
months earlier.

\begin{dated}{2026-09}{Linker markets}\label{dat:m2:inflation-markets:markets}
\textbf{United States}: marketable TIPS outstanding USD~2\,152.7 billion on 31
August 2026, 6.8\% of USD~31\,828.0 billion of marketable Treasury debt; terms
of 5, 10 and 30 years. \textbf{United Kingdom}: the Retail Prices Index is to
be calculated with the methods and data of the consumer price index including
owner occupiers' housing costs (CPIH) from February 2030, as the Statistics
Authority and the Treasury announced on 25 November 2020, a date chosen to
limit the effect on holders of index-linked gilts. \textbf{Euro area}:
Eurostat rebased the harmonised index to 2025 from February 2026; the French Treasury states that the rebasing does not
change the indexation coefficients of its bonds.
\end{dated}

\section{Real yields and breakevens}

\begin{definition}[Real yield, breakeven inflation rate]\label{def:m2:inflation-markets:breakeven}
The \emph{real yield}\index{real yield} of an \omterm{def:m2:inflation-markets:linker}{index-linked bond} is the yield to
maturity of its real cash flows (the fixed coupon and the face value, in units
of today's index) at its real price, the quoted price per 100 of adjusted
principal. The \emph{breakeven inflation rate}\index{breakeven inflation rate}
at a maturity is the constant inflation rate at which a nominal bond and a \omterm{def:m2:inflation-markets:linker}{linker}
of that maturity have the same return:
\begin{equation}\label{eq:m2:inflation-markets:fisher}
1 + y_{\mathrm{nom}} \;=\; (1 + y_{\mathrm{real}})(1 + \pi_{\mathrm{BE}}),
\qquad \pi_{\mathrm{BE}} \approx y_{\mathrm{nom}} - y_{\mathrm{real}}.
\end{equation}
\end{definition}

On 29 July 2022 the ten-year nominal Treasury yield was 2.67\% and the ten-year
\omterm{def:m2:inflation-markets:breakeven}{real yield} 0.14\%: the exact breakeven is $1.0267/1.0014 - 1 = 2.5265\%$ and the
difference 2.53\%, the figure screens show. An investor who expects inflation
above the breakeven prefers the \omterm{def:m2:inflation-markets:linker}{linker}; one who expects less, the nominal bond.

\begin{proposition}[What a breakeven contains]\label{prop:m2:inflation-markets:contents}
Under the risk-neutral measure the breakeven equals expected inflation at that
horizon only if investors demand no premium for bearing inflation risk and
the two bonds are equally liquid. In general
\[
\pi_{\mathrm{BE}} \;\approx\; \mathbb{E}[\pi] + \text{inflation risk premium}
- \text{liquidity premium of the linker} + \text{convexity terms},
\]
with the liquidity term entering negatively because a less liquid \omterm{def:m2:inflation-markets:linker}{linker} must
offer a higher \omterm{def:m2:inflation-markets:breakeven}{real yield}, which lowers the measured breakeven.
\end{proposition}
\admitted{The decomposition is a statement about the pricing kernel; term-structure
models that estimate each part are the subject of One Quant Book 6.}

The practical consequence: a breakeven is a price, not a forecast. It falls
when investors flee to the most liquid assets, as in March 2020, when the
ten-year breakeven fell to 0.50\% on 19 March, and it can stay below realised
inflation for long stretches.

\begin{omfigure}
\begin{tikzpicture}
\begin{axis}[width=11cm,height=5.2cm,
  xlabel={},ylabel={percent},
  tick label style={font=\small},label style={font=\small},
  xmin=2018,xmax=2026.75,ymin=-1.5,ymax=5.5,grid=major,grid style={gray!25},
  xtick={2018,2019,2020,2021,2022,2023,2024,2025,2026},xticklabel style={/pgf/number format/1000 sep={}},
  legend columns=3,legend style={font=\footnotesize,at={(0.5,-0.18)},anchor=north,draw=none,fill=none,/tikz/every even column/.append style={column sep=8pt}},
  table/col sep=comma]
\addplot[omDef,very thick] table[x=t,y=nominal]{figdata/markets-2/11-inflation-markets/tips10.csv};
\addplot[omThm,very thick] table[x=t,y=real]{figdata/markets-2/11-inflation-markets/tips10.csv};
\addplot[omProp,thick,dashed] table[x=t,y=breakeven]{figdata/markets-2/11-inflation-markets/tips10.csv};
\draw[gray] (axis cs:2018,0) -- (axis cs:2026.75,0);
\legend{10-year nominal,10-year real (TIPS),10-year breakeven}
\end{axis}
\end{tikzpicture}

\omcaption{Ten-year nominal and real Treasury yields and the breakeven between
them, month ends, January 2018 to August 2026. In 2022 the \omterm{def:m2:inflation-markets:breakeven}{real yield} rose by
more than two and a half percentage points while the breakeven hardly moved:
the losses of inflation-linked bonds that year came from real rates. Data: FRED
series DGS10, DFII10 and T10YIE.}\label{fig:m2:inflation-markets:tips10}
\end{omfigure}

The figure also answers the hook. From the start to the end of 2022 the
breakeven fell a little, from 2.60\% to 2.30\%, while the \omterm{def:m2:inflation-markets:breakeven}{real yield} rose by 2.55
percentage points. A \omterm{def:m2:inflation-markets:linker}{linker}'s price moves with its \omterm{def:m2:inflation-markets:breakeven}{real yield} exactly as a
nominal bond's moves with its nominal yield, with a real duration of the same
size; the indexation is added on top. The ten-year TIPS of the hook had a real
duration above nine years, so the move cost it about a fifth of its value,
while a year of 9\% inflation, reaching the principal with a lag, added 7.7\%.

\section{Zero-coupon inflation swaps}

\begin{definition}[Zero-coupon inflation swap]\label{def:m2:inflation-markets:zciis}
A \emph{zero-coupon inflation swap}\index{zero-coupon inflation swap} of notional
$N$, maturity $T$ years and fixed rate $K$ exchanges a single net payment at
maturity: the inflation receiver gets
\[
N\left(\frac{I(T)}{I(0)} - (1+K)^T\right),
\]
where $I(0)$ and $I(T)$ are the \omterm{def:m2:inflation-markets:linker}{reference index} at the start and at maturity,
with the same lag as the bonds. Nothing is exchanged at the start.
\end{definition}

Dollar inflation swaps reference the same non-seasonally-adjusted index as
TIPS, and the zero-coupon swap is the most common form: a study of dealers'
trades by the Federal Reserve Bank of New York counted 144 trades over June to
August 2010, about USD~65 million a day, against about USD~5 billion a day of
TIPS. Pension funds and insurers, whose liabilities grow with prices, receive
inflation; issuers whose revenues are indexed, such as utilities, pay it.

\begin{proposition}[The forward index]\label{prop:m2:inflation-markets:forward}
Since the swap is worth zero at inception, the market's zero-coupon rate $K_T$
fixes a forward value of the index: $I^{\mathrm{fwd}}(T) = I(0)(1+K_T)^T$ is the
level at which the inflation leg and the fixed leg are equal, and any
index-linked cash flow $c\,I(T)/I(0)$ paid at $T$ is worth
$c(1+K_T)^T P(0,T)$ today, where $P(0,T)$ is the nominal discount factor.
\end{proposition}

\begin{proof}
Receive the index-linked flow and pay inflation on a swap of notional $c$:
the index terms cancel and $c(1+K_T)^T$, known today, is left at $T$. Its
value is $c(1+K_T)^T P(0,T)$; the swap was free.
\end{proof}

\begin{example}[A five-year swap]\label{ex:m2:inflation-markets:swap}
A pension fund receives inflation on USD~100 million for five years at 2.50\%,
with a base index of 300.00. The forward index is $300 \times 1.025^5 = 339.42$.
If the index ends at 345.00, 15\% above its base, the fund receives
$100 \times (1.15 - 1.025^5) = \text{USD}~1.86$ million; had it ended at the
forward, nothing.
\end{example}

The proposition turns a curve of swap rates into a curve of forward index
values, and so into a price for every \omterm{def:m2:inflation-markets:linker}{linker}: discount its indexed cash flows
at the nominal curve. The gap between that price and the market price, quoted
as a spread over the swap curve, moves with the balance-sheet cost of holding
bonds, as the \omterm{def:m2:interest-rate-swaps:spread}{swap spread} of \cref{ch:m2:interest-rate-swaps} does
(\cref{ch:m2:corporate-bonds} returns to such spreads).

\section{Indexation lag and seasonality}

\begin{definition}[Indexation lag, inflation seasonality]\label{def:m2:inflation-markets:lag}
The \emph{indexation lag}\index{indexation lag} of a \omterm{def:m2:inflation-markets:linker}{linker} is the delay between
the month whose prices set its \omterm{def:m2:inflation-markets:linker}{reference index} and the month in which that
reference applies: three months for TIPS and for most bonds issued since the
2000s. \emph{Inflation seasonality}\index{inflation seasonality} is the recurring
pattern of a non-seasonally-adjusted price index within the calendar year,
from the timing of sales, energy use and annual price resets.
\end{definition}

Two consequences follow, and both matter to a trader. The first is in
\cref{fig:m2:inflation-markets:lag}: the growth of the \omterm{def:m2:inflation-markets:linker}{reference index} over a
month is the growth of the published index two and three months earlier,
\begin{equation}\label{eq:m2:inflation-markets:accrual}
\frac{\mathrm{Ref}(\text{1st of } M+1)}{\mathrm{Ref}(\text{1st of } M)} - 1 \;=\; \frac{I_{M-2}}{I_{M-3}} - 1,
\end{equation}
so a \omterm{def:m2:inflation-markets:linker}{linker}'s inflation accrual for the next month or two is known exactly. In
2022 the months of the largest prints followed each other: the \omterm{def:m2:inflation-markets:linker}{reference index}
grew 1.37\% over August alone, 17.8\% at an annual rate, then fell slightly in
September (\cref{fig:m2:inflation-markets:accrual}).

\begin{omfigure}
\begin{tikzpicture}
\begin{axis}[width=11cm,height=4.8cm,ybar,bar width=4pt,
  ylabel={accrual over the month (\%)},
  tick label style={font=\small},label style={font=\small},
  xmin=-1,xmax=36,ymin=-0.5,ymax=1.95,ytick={-0.5,0,0.5,1,1.5},grid=major,grid style={gray!25},
  xtick={0,12,24},xticklabels={Jan 2021,Jan 2022,Jan 2023},
  table/col sep=comma]
\addplot[fill=omDef!70,draw=omDef] table[x=k,y=accrual]{figdata/markets-2/11-inflation-markets/accrual.csv};
\draw[omThm,thick,-{Stealth}] (axis cs:17.3,1.72) -- (axis cs:18.8,1.45);
\node[omThm,font=\footnotesize,anchor=east,align=right] at (axis cs:17.3,1.66) {August 2022: 1.37\%,\\known on 13 July};
\end{axis}
\end{tikzpicture}

\omcaption{Growth of the TIPS \omterm{def:m2:inflation-markets:linker}{reference index} over each calendar month, January
2021 to December 2023: the monthly change in the non-seasonally-adjusted index
of three and two months earlier. Data: FRED series CPIAUCNS, the chapter's
tutorial.}\label{fig:m2:inflation-markets:accrual}
\end{omfigure}

The second consequence is seasonality. Because the index is not seasonally
adjusted, the monthly accrual has a calendar pattern on top of the trend:
prices tend to rise faster early in the year and slower, or to fall, late in it
(\cref{fig:m2:inflation-markets:seasonal}). Over a ten-year bond the pattern
averages out. Over a few months it does not, and a short-dated breakeven read
without adjusting for it gives a wrong signal: a bond whose remaining accrual
covers the reference months of October to December carries about 0.85 percentage
points less inflation than the average quarter, 3.4 points a year annualised.

\begin{omfigure}
\begin{tikzpicture}
\begin{axis}[width=11cm,height=4.6cm,ybar,bar width=12pt,
  ylabel={excess change (pp)},xlabel={month of the index},
  tick label style={font=\small},label style={font=\small},
  xmin=0.3,xmax=12.7,ymin=-0.45,ymax=0.4,grid=major,grid style={gray!25},
  xtick={1,2,3,4,5,6,7,8,9,10,11,12},xticklabels={Jan,Feb,Mar,Apr,May,Jun,Jul,Aug,Sep,Oct,Nov,Dec},
  table/col sep=comma]
\addplot[fill=omThm!60,draw=omThm] table[x=month,y=factor]{figdata/markets-2/11-inflation-markets/seasonal.csv};
\end{axis}
\end{tikzpicture}

\omcaption{Seasonality of the US consumer price index, not seasonally adjusted:
the average change of each month's index over the previous month, minus the
average monthly change of its year, 2010 to 2024 (2025 lacks its October
index), in percentage points. March
is the strongest month, November and December the weakest. Data: FRED series
CPIAUCNS, the chapter's tutorial.}\label{fig:m2:inflation-markets:seasonal}
\end{omfigure}

\begin{proposition}[Carry of a financed linker]\label{prop:m2:inflation-markets:carry}
A \omterm{def:m2:inflation-markets:linker}{linker} of market value $V$ held for $\tau$ days with an unchanged \omterm{def:m2:inflation-markets:breakeven}{real yield}
$y$, financed in repo at rate $r$, earns
\[
\mathrm{carry} \;=\; V\left[(1+a)(1+y)^{\tau/365} - 1\right] - V r\,\frac{\tau}{360},
\]
where $a$ is the growth of the \omterm{def:m2:inflation-markets:linker}{reference index} over the period, known in
advance for up to about two months.
\end{proposition}

\begin{proof}
With the \omterm{def:m2:inflation-markets:breakeven}{real yield} unchanged the real price grows at the \omterm{def:m2:inflation-markets:breakeven}{real yield} (the
pull to par plus the coupon accrual), and the \omterm{def:m2:bond-futures:cf}{invoice price} is the real price
times the index ratio, which grows by $1+a$. Repo interest is paid on $V$ on an
actual/360 basis.
\end{proof}

Because $a$ is known, it is in the price: a forward sale of the \omterm{def:m2:inflation-markets:linker}{linker} for the
end of the month is struck at a price that includes it, and a buyer ahead of a
large print pays for it. What remains uncertain, and what a trader earns or
loses, is the move in the \omterm{def:m2:inflation-markets:breakeven}{real yield} and the prints not yet published.

\section{Tutorial: breakevens, seasonality and carry}\label{tut:m2:inflation-markets:carry}

\begin{tutorial}
\textbf{Goal.} Compute reference indices and index ratios by the Treasury's
rule, estimate the seasonal pattern of the index, and price the carry of a
financed TIPS. \textbf{End state:}
\cref{fig:m2:inflation-markets:tips10,fig:m2:inflation-markets:accrual,fig:m2:inflation-markets:seasonal},
\cref{ex:m2:inflation-markets:ratio} and the numbers of the weekend problem.
\begin{tutsteps}
\item \textbf{The \omterm{def:m2:inflation-markets:linker}{reference index}}, the index ratio and the month's accrual,
by \cref{eq:m2:inflation-markets:refcpi,eq:m2:inflation-markets:accrual}.
\omcode{code/firm/breakeven/firm_breakeven.py}{26}{41}{Reference index, index ratio and monthly accrual.}\label{lst:m2:inflation-markets:ref}
\item \textbf{Seasonality and carry.} Each month's log change less its year's
mean, averaged over complete years (a year with a missing month is skipped); the carry of
\cref{prop:m2:inflation-markets:carry}.
\omcode{code/firm/breakeven/firm_breakeven.py}{59}{81}{Seasonal factors and the carry of a financed linker.}\label{lst:m2:inflation-markets:carry}
\item \textbf{Run} \texttt{inflation\_demo.hook\_2022()},
\texttt{inflation\_demo.august\_2022()} and \texttt{fig\_inflation.py}, which
read \texttt{data/markets-2/cpi\_us\_monthly.csv} and
\texttt{tips10\_monthend.csv}.
\end{tutsteps}
\textbf{What to change next.} Estimate the seasonal factors on 2010--2019 only
and compare; then compute the carry of September 2022 and explain its sign.
\end{tutorial}

\section{Build: the breakeven and carry calculator}\label{bld:m2:inflation-markets:breakeven}

\begin{build}
\bfield{Purpose} The miniature firm quotes and finances inflation-linked bonds
and swaps: it needs the \omterm{def:m2:inflation-markets:linker}{reference index} of any day, each bond's index ratio,
breakevens, forward index values from swap rates, and the carry it will earn
over the next month.
\bfield{Interface} \texttt{ref\_index(d, index, lag)};
\texttt{index\_ratio(d, dated, index, lag)}; \texttt{monthly\_accrual(year,
month, index, lag)}; \texttt{breakeven(nominal, real)};
\texttt{forward\_index(base, zc\_rate, years)}; \texttt{zc\_swap\_payment};
\texttt{seasonal\_factors(index)}; \texttt{linker\_carry(value, accrual,
real\_yield, repo, days)}; \texttt{contingency\_index(last,
year\_before\_last, missing)}.
\bfield{Rules} Index values keyed by (year, month); \omterm{def:m2:inflation-markets:linker}{reference index} and index
ratio truncated to six decimals and rounded to five; a lag parameter so that
eight-month gilts are priced by the same code; seasonal factors only from
complete calendar years.
\bfield{Acceptance tests} \texttt{code/firm/breakeven/tests/}: the regulation's
own example (154.63333 and 1.00011); the first of a month equals the lagged
index; the October 2025 substitute (325.604); a swap struck at the forward pays nothing; seasonal factors sum to zero;
carry signs.
\bfield{Stretch} Seasonally adjusted forward index curves from swap rates; a
real-yield solver for a TIPS from its quoted real price; the UK eight-month
lag gilts, whose coupons are fixed in advance.
\end{build}

\begin{omsources}
\item Code of Federal Regulations, 31 CFR Part 356, Appendix B (formulas for
inflation-protected securities); TreasuryDirect, TIPS and the history of TIPS.
\item UK Debt Management Office, FAQs on index-linked gilts; UK Statistics
Authority, response to the joint consultation on reforming the Retail Prices
Index, 25 November 2020.
\item Agence France Tr\'esor, characteristics of OATi and OAT\texteuro i.
\item M.~J.~Fleming and J.~R.~Sporn, ``Trading activity and price
transparency in the inflation swap market'', \emph{Federal Reserve Bank of New
York Economic Policy Review}, May 2013.
\item Bureau of Labor Statistics, Consumer Price Index news release for June
2022; FRED series CPIAUCNS, CPIAUCSL, DGS10, DFII10, T10YIE; Treasury, Monthly
Statement of the Public Debt.
\end{omsources}

\section{Exercises}

\begin{exercise}[$\star$]\label{exo:m2:inflation-markets:1}
With the May 2022 index at 292.296 and June at 296.311, give the TIPS
\omterm{def:m2:inflation-markets:linker}{reference index} on 15 August 2022.
\end{exercise}

\begin{exercise}[$\star$]\label{exo:m2:inflation-markets:2}
The ten-year nominal yield is 2.67\% and the ten-year \omterm{def:m2:inflation-markets:breakeven}{real yield} 0.14\%. Give
the exact breakeven and the approximation.
\end{exercise}

\begin{exercise}[$\star$]\label{exo:m2:inflation-markets:3}
Over the life of a bond the index falls 3\%. What does a TIPS repay per 100 of
face value, and what does an index-linked gilt repay?
\end{exercise}

\begin{exercise}[$\star\star$]\label{exo:m2:inflation-markets:4}
A fund receives inflation on USD~100 million for five years at 2.50\%, base
index 300.00. The index ends at 345.00. What is the net payment, and at what
final index would there be none?
\end{exercise}

\begin{exercise}[$\star\star$]\label{exo:m2:inflation-markets:5}
Give the index ratio on 15 August 2022 of the TIPS of
\cref{ex:m2:inflation-markets:ratio}, and explain why the coupon paid on 15
July used a smaller ratio.
\end{exercise}

\begin{exercise}[$\star\star$]\label{exo:m2:inflation-markets:6}
With the factors of \cref{fig:m2:inflation-markets:seasonal}, by how much does
seasonality lower the inflation a TIPS accrues over the three months whose
reference indices are those of October, November and December, and by how much
at an annual rate?
\end{exercise}

\begin{exercise}[$\star\star\star$]\label{exo:m2:inflation-markets:7}
\emph{Coding.} Estimate the seasonal factors on 2010--2019 only. Which months
are the strongest and the weakest, and how far does any month move from the
2010--2024 estimate?
\end{exercise}

\begin{exercise}[$\star\star\star$]\label{exo:m2:inflation-markets:8}
\emph{Find the flaw.} ``Inflation is running at 9\% and the ten-year breakeven
is 2.5\%: buy TIPS, sell nominal Treasuries, and collect the difference.''
Correct it.
\end{exercise}

\section{Problem: The Linker's Carry}

\begin{problem}[{Weekend problem --- what a TIPS earns in a month of a big print}]\label{pb:m2:inflation-markets:1}
On 13 July 2022 the June index was published at 296.311, 1.37\% above May's
292.296. A trader holds USD~100 million face of the ten-year TIPS of
\cref{ex:m2:inflation-markets:ratio} (coupon 0.125\%, dated 15 January 2022)
and finances it in repo at 2.30\%, about the overnight rate in early August. On
29 July the ten-year \omterm{def:m2:inflation-markets:breakeven}{real yield} was 0.14\% and the nominal yield 2.67\%;
assume they do not move during August.

\medskip\noindent\textbf{Part I --- The index.}
\begin{enumerate}
\item Give the \omterm{def:m2:inflation-markets:linker}{reference index} on 1 August and on 1 September 2022.
\item Give the index ratio on 1 August.
\item Give the growth of the \omterm{def:m2:inflation-markets:linker}{reference index} over August.
\item When did the trader know it, and why?
\item Give the \omterm{def:m2:inflation-markets:linker}{reference index} on 15 August.
\end{enumerate}

\medskip\noindent\textbf{Part II --- The carry.}
\begin{enumerate}[resume]
\item Give the real price at 0.14\% and the position's market value on 1 August.
\item Give the income from the index accrual over August's 31 days.
\item Give the income from the \omterm{def:m2:inflation-markets:breakeven}{real yield}.
\item Give the repo cost.
\item Give the carry.
\end{enumerate}

\medskip\noindent\textbf{Part III --- Comparisons.}
\begin{enumerate}[resume]
\item Give the carry of the same market value of a nominal bond yielding 2.67\%.
\item Give August's accrual at an annual rate, and compare it with the
breakeven.
\item The July index was 296.276. Give September's accrual and carry (30 days).
\item Give the position's DV01 in \omterm{def:m2:inflation-markets:breakeven}{real yield}, and the real-yield move that
would cancel August's carry.
\item Why is the carry larger for a short-dated TIPS, relative to its risk?
\end{enumerate}

\medskip\noindent\textbf{Part IV --- Judgement.}
\begin{enumerate}[resume]
\item Is the carry a profit for a trader who buys on 29 July?
\item Why did TIPS lose money in 2022 despite such months?
\item What is the \omterm{def:m2:inflation-markets:linker}{deflation floor} of this bond worth in August 2022?
\item State the \emph{named result}: the carry of USD~100 million face of the
ten-year TIPS over August 2022.
\item In one sentence: what does a \omterm{def:m2:inflation-markets:linker}{linker} protect against, and what not?
\end{enumerate}
\end{problem}

\section{Interview questions}

\begin{interviewq}[$\star$ \iqroles{trader, researcher}]\label{iq:m2:inflation-markets:1}
What is a \omterm{def:m2:inflation-markets:breakeven}{breakeven inflation rate}, and why is it not the market's inflation
forecast?
\end{interviewq}

\begin{interviewq}[$\star$ \iqroles{bank, developer}]\label{iq:m2:inflation-markets:2}
How does the principal of a TIPS change, and what happens at maturity after
deflation?
\end{interviewq}

\begin{interviewq}[$\star\star$ \iqroles{trader}]\label{iq:m2:inflation-markets:3}
Why does a TIPS carry so much in a month after a large inflation print, and
who gets that carry?
\end{interviewq}

\begin{interviewq}[$\star\star$ \iqroles{researcher, trader}]\label{iq:m2:inflation-markets:4}
The five-year zero-coupon inflation swap rate and the five-year TIPS breakeven
differ. What could make them differ?
\end{interviewq}

\begin{interviewq}[$\star\star$ \iqroles{researcher}]\label{iq:m2:inflation-markets:5}
How would you remove seasonality from a price index to price short-dated
\omterm{def:m2:inflation-markets:linker}{linkers}?
\end{interviewq}

\begin{interviewq}[$\star\star\star$ \iqroles{developer, bank}]\label{iq:m2:inflation-markets:6}
Design the service that publishes the index ratio of every \omterm{def:m2:inflation-markets:linker}{linker} a firm holds,
in the United States, the United Kingdom and the euro area, every morning.
\end{interviewq}
